\SourceLeaf{167}ja Deine Arbeit bald bekannt; gewiss wirst Du dafür den Dank nicht zwar
des grossen Publikums (worunter auch mancher gehört, der für einen
geschickten Mathematiker gehalten wird) einernten, denn ich überzeuge
mich immer mehr, dass die Zahl wahrer Geometer äusserst gering ist und
die meisten die Schwierigkeiten bei solchen Arbeiten weder beurtheilen
noch selbst einmal sie verstehen können --- aber gewiss den Dank aller
derer, deren Urtheil Dir allein wirklich schätzbar sein kann. --- In
Braunschweig ist ein Emigrant Namens \textsc{Chauvelot}, ein nicht
schlechter Geometer, welcher vorgibt, die Theorie der Parallellinien
ganz begründet zu haben und seine Arbeit nächstens wird drucken
lassen, aber ich verspreche mir eben nichts von ihm. In
\textsc{Hindenburgs} Archiv, 9\textsuperscript{tes} Stück, befindet
sich gleichfalls ein neuer Versuch über denselben Gegenstand, von
einem gewissen \textsc{Hauff}, welcher unter aller Kritik ist. .\, .\, .\, .\, .\, .\, .

\medskip
\begin{center}
\rule{0.3\textwidth}{0.4pt}
\end{center}
\medskip

\begin{center}
\textsc{Gauss} an \textsc{Bolyai}. Braunschweig, 25. November 1804.
\end{center}

.\, .\, .\, .\, .\, .\, .\ Ich habe Deinen Aufsatz mit grossem Interesse und
Aufmerksamkeit durchgelesen, und mich recht an dem ächten gründlichen
Scharfsinne ergötzt. Du willst aber nicht mein leeres Lob, das auch
gewissermassen schon darum parteiisch scheinen könnte, weil Dein
Ideengang sehr viel mit dem meinigen Ähnliches hat, worauf ich
ehemals die Lösung dieses Gordischen Knoten versuchte, und vergebens
bis jetzt versuchte. Du willst nur mein aufrichtiges unverholenes
Urtheil. Und diess ist, dass Dein Verfahren mir noch
\textls{nicht} 
Genüge leistet. Ich will versuchen, den Stein des Anstosses, den ich
noch darin finde (und der auch wieder zu derselben Gruppe von Klippen
gehört, woran meine Versuche bisher scheiterten) mit so vieler
Klarheit, als mir möglich ist, ans Licht zu ziehen. Ich habe zwar noch
immer die Hoffnung, dass jene Klippen einst, und noch vor meinem
Ende, eine Durchfahrt erlauben werden. Indess habe ich jetzt so
manche andere Beschäftigungen vor der Hand, dass ich gegenwärtig
daran nicht denken kann, und glaube mir, es soll mich herzlich
herzlich freuen, wenn Du mir zuvorkommst, und es Dir
